\documentclass[11pt,a4paper]{article}
\usepackage[utf8]{inputenc}
\usepackage[T1]{fontenc}
\usepackage[polish]{babel}
\usepackage{lmodern}
\usepackage[margin=2.0cm]{geometry}
\usepackage{booktabs}
\usepackage{tabularx}
\usepackage{array}
\usepackage{enumitem}
\usepackage{titlesec}
\usepackage{xcolor}
\usepackage{fancyhdr}
\usepackage{hyperref}
\usepackage{url}

\definecolor{navy}{RGB}{27,42,86}
\definecolor{lgray}{RGB}{240,240,244}

\titleformat{\section}{\large\bfseries\color{navy}}{\thesection.}{0.5em}{}
\titleformat{\subsection}{\bfseries\color{navy}}{\thesubsection.}{0.5em}{}

\pagestyle{fancy}
\fancyhf{}
\renewcommand{\headrulewidth}{0.6pt}
\fancyhead[L]{\small Neurofizjologia i pomiar EEG}
\fancyhead[R]{\small Instrukcja laboratoryjna nr 1}
\fancyfoot[C]{\small\thepage}

\newcommand{\infobox}[1]{%
  \par\noindent\colorbox{lgray}{\parbox{\dimexpr\linewidth-2\fboxsep}{\vspace{2pt}#1\vspace{2pt}}}\par\vspace{4pt}}

\title{\vspace{-1.5cm}\textbf{\color{navy}Instrukcja laboratoryjna nr 1}\\[4pt]
\Large Wprowadzenie do badań elektrofizjologicznych}
\author{\vspace{-1.5cm}dr inż. Mateusz Pomianek}
\date{\vspace{-1.4cm}\today}


\begin{document}
\maketitle
\thispagestyle{fancy}

\section{Cele zajęć}
\begin{itemize}[nosep]
  \item Zapoznanie studentów z historycznym i teoretycznym podłożem elektroencefalografii (EEG).
  \item Wyjaśnienie neurofizjologicznych źródeł sygnału EEG i jego związku z aktywnością korową.
  \item Przegląd głównych zastosowań EEG w badaniach nad zaburzeniami psychicznymi i funkcjonowaniem poznawczym.
  \item Wprowadzenie do podstawowej terminologii wykorzystywanej w dalszych zajęciach laboratoryjnych.
\end{itemize}

\section{Wprowadzenie teoretyczne}

\subsection{Rys historyczny}
Hans Berger w 1929 roku jako pierwszy zarejestrował i opisał ludzką elektroencefalografię, odkrywając rytm alfa. Wcześniejsze prace Richarda Catona (1875) na zwierzętach wykazały istnienie prądów elektrycznych w mózgu. Rozwój wzmacniaczy elektronicznych w latach 30.--50. XX wieku oraz komputeryzacja analizy sygnału od lat 70. umożliwiły przejście od klinicznej obserwacji zapisu papierowego do ilościowej analizy EEG (qEEG).

\subsection{Źródła sygnału EEG}
Sygnał rejestrowany na powierzchni skóry głowy jest sumą \textbf{postsynaptycznych potencjałów wzbudzających i hamujących (EPSP/IPSP)} generowanych głównie przez neurony piramidowe warstw III i V kory mózgowej, których dendryty ułożone są prostopadle do powierzchni kory (organizacja \emph{open field}). Pojedynczy potencjał czynnościowy jest zbyt krótki i słaby, by wygenerować mierzalny sygnał powierzchniowy -- konieczna jest \textbf{synchroniczna aktywacja} dużej liczby (rzędu $10^4$--$10^7$) sąsiadujących neuronów.

\subsection{Pasma częstotliwości}
\begin{center}
\begin{tabular}{@{}lll@{}}
\toprule
\textbf{Pasmo} &
\textbf{Zakres (Hz)} &
\textbf{Typowy kontekst wystąpienia}\\
\midrule
Delta ($\delta$) & 0,5--4 & Sen głęboki (NREM N3), patologia u dorosłych w czuwaniu\\
Theta ($\theta$) & 4--8 & Senność, procesy pamięciowe, dzieci\\
Alfa ($\alpha$) & 8--13 & Czuwanie z zamkniętymi oczami, relaksacja\\
Beta ($\beta$) & 13--30 & Czuwanie aktywne, koncentracja, napięcie\\
Gamma ($\gamma$) & $>$30 & Wiązanie percepcyjne, procesy poznawcze wyższego rzędu\\
\bottomrule
\end{tabular}
\end{center}

\subsection{Strukturalne i funkcjonalne obszary mózgu a topografia zapisu EEG}
Rozmieszczenie elektrod w systemie 10--20 (szczegółowo omawiane na zajęciach 2) nie jest przypadkowe -- odpowiada w przybliżeniu położeniu głównych płatów kory leżących pod danym obszarem czaszki. Znajomość tej zgrubnej korespondencji jest niezbędna do sensownej interpretacji topografii sygnału.

\begin{center}
\begin{tabular}{p{0.1\textwidth} p{0.35\textwidth} p{0.45\textwidth}}
\toprule
\textbf{Elektrody} &
\textbf{Leżąca struktura} &
\textbf{Powiązana funkcja / aktywność} \\
\midrule
Fp, F & Płat czołowy (kora przedczołowa, kora ruchowa dodatkowa) & Funkcje wykonawcze, planowanie, kontrola uwagi; asymetria pasma alfa jako korelat regulacji emocji\\
C & Płat ciemieniowy przedśrodkowy / kora ruchowa i czuciowa pierwotna & Planowanie i wykonywanie ruchu; rytm mu (8--13 Hz), tłumiony podczas ruchu i jego obserwacji\\
T & Płat skroniowy (kora słuchowa, struktury przyśrodkowo-skroniowe) & Przetwarzanie słuchowe i językowe, procesy pamięciowe (aktywność theta)\\
P & Płat ciemieniowy & Integracja czuciowo-przestrzenna, uwaga przestrzenna\\
O & Płat potyliczny (kora wzrokowa pierwotna i asocjacyjna) & Przetwarzanie wzrokowe; główny generator dominującego rytmu alfa przy zamkniętych oczach\\
\bottomrule
\end{tabular}
\end{center}

\textbf{Ważne zastrzeżenie -- przewodnictwo objętościowe.} Sygnał elektryczny generowany w korze rozprzestrzenia się przez tkanki o różnej przewodności (opony mózgowe, płyn mózgowo-rdzeniowy, czaszka, skóra) zanim dotrze do elektrody -- zjawisko to nazywa się \emph{przewodnictwem objętościowym (volume conduction)}. W efekcie potencjał zarejestrowany w danym punkcie skalpu jest mieszaniną aktywności wielu, czasem odległych, generatorów korowych, a nie \guillemotleft fotografią\guillemotright{} leżącej bezpośrednio pod nim kory. Ponadto EEG jest najbardziej czułe na źródła o orientacji dipolowej zbliżonej do prostopadłej względem powierzchni czaszki (typowe dla neuronów w grzbietach zakrętów), słabiej rejestrując aktywność generowaną w głębi bruzd. Zagadnienie to jest bezpośrednio powiązane z problemem analizy źródłowej, omawianym szczegółowo na zajęciach 4.

\subsection{EEG w badaniach zaburzeń psychicznych i funkcjonowania poznawczego}
\begin{center}
\begin{tabular}{p{0.20\textwidth} p{0.70\textwidth}}
\toprule
\textbf{Obszar kliniczny} &
\textbf{Typowe nieprawidłowości / zastosowania EEG}\\
\midrule
Schizofrenia & Redukcja amplitudy komponentu P300, zaburzenia synchronizacji oscylacji gamma (deficyt \emph{binding}), nadmiar wolnofalowej aktywności theta/delta\\
Zaburzenia nastroju & Asymetria czołowa pasma alfa (relatywnie niższa aktywacja lewej kory przedczołowej) jako korelat podatności na depresję\\
ADHD & Podwyższony (choć dyskusyjny jako samodzielny biomarker) stosunek mocy theta/beta w odprowadzeniach czołowo-środkowych\\
Otępienia & Postępujące spowolnienie rytmu podstawowego (przesunięcie w kierunku theta/delta), zmniejszenie zmienności i złożoności sygnału\\
Padaczka & Interiktalne wyładowania iglica--fala (spike-wave), lokalizacja ogniska padaczkorodnego\\
Zaburzenia snu & Polisomnografia -- klasyfikacja stadiów snu na podstawie wzorców EEG, EOG i EMG\\
Funkcjonowanie poznawcze & Potencjały wywołane (P300, N400) i moc pasmowa jako korelaty uwagi, pamięci roboczej i przetwarzania językowego\\
\bottomrule
\end{tabular}
\end{center}

\infobox{\textbf{Do zapamiętania:} EEG mierzy zsumowaną aktywność synaptyczną, nie pojedyncze potencjały czynnościowe. Rozdzielczość czasowa jest rzędu milisekund, natomiast przestrzenna -- ograniczona przez przewodnictwo objętościowe do rzędu centymetrów.}

\section{Sprzęt i oprogramowanie wykorzystywane na zajęciach}
Prezentacja poglądowa: wzmacniacz EEG, przykładowy czepek z elektrodami, oprogramowanie do akwizycji i wizualizacji sygnału (np. BrainVision Recorder/Analyzer, EEGLAB w środowisku MATLAB lub Python/MNE). Szczegółowe omówienie budowy i zasady działania aparatury oraz procedury jej przygotowania stanowi przedmiot zajęć 2.

\section{Przebieg zajęć}
\begin{center}
\begin{tabular}{p{0.15\textwidth} p{0.80\textwidth}}
\toprule
\textbf{Czas} &
\textbf{Aktywność}\\
\midrule
0--10 min & Wprowadzenie, omówienie celów zajęć i sylabusa laboratorium.\\
10--30 min & Wykład: historia EEG, neurofizjologiczne źródła sygnału, pasma częstotliwości.\\
30--50 min & Wykład: strukturalne i funkcjonalne obszary mózgu a topografia zapisu EEG, przewodnictwo objętościowe.\\
50--60 min & Przerwa.\\
60--80 min & Wykład: zastosowania EEG w zaburzeniach psychicznych i funkcjonowaniu poznawczym.\\
80--105 min & Praca w grupach: analiza fragmentów opublikowanych badań EEG (zadanie 1) i mapowanie płat--funkcja--pasmo (zadanie 2).\\
105--115 min & Krótka prezentacja wniosków grup i dyskusja.\\
115--120 min & Podsumowanie, zapowiedź kolejnych zajęć (sprzęt i procedura badania).\\
\bottomrule
\end{tabular}
\end{center}

\section{Zadania do wykonania}
\begin{enumerate}[label=\textbf{Zadanie \arabic*.}, leftmargin=2.7cm]
  \item W grupach 3--4-osobowych przeanalizujcie przydzielony przez prowadzącego artykuł naukowy wykorzystujący EEG (np. w kontekście ADHD, schizofrenii lub uwagi). Zidentyfikujcie: cel badania, badane pasmo/miary, wnioski kliniczne.
  \item Na schemacie rozmieszczenia elektrod 10--20 zaznaczcie i podpiszcie granice czterech głównych płatów mózgu, a dla każdego obszaru wskażcie jedną typową funkcję poznawczą oraz jedno dominujące pasmo/zjawisko EEG z nim związane.
  \item Przygotujcie 3-minutową prezentację ustną podsumowującą analizowany artykuł (zadanie 1), ze wskazaniem mocnych i słabych stron metodologicznych.
  \item Indywidualnie: sporządźcie tabelę porównawczą pięciu pasm częstotliwości EEG wraz z przykładami zastosowań klinicznych (do oddania na kolejnych zajęciach).
\end{enumerate}

\section{Kryteria zaliczenia ćwiczenia}
Aktywny udział w dyskusji grupowej i poprawność mapowania płat--funkcja--pasmo (40\%), poprawność merytoryczna prezentacji artykułu (30\%), oddana tabela porównawcza pasm (30\%).

\section{Pytania kontrolne}
\begin{enumerate}[nosep]
  \item Dlaczego pojedynczy potencjał czynnościowy neuronu nie jest widoczny w zapisie EEG?
  \item Jakie znaczenie ma organizacja \emph{open field} neuronów piramidowych dla powstawania sygnału EEG?
  \item Na czym polega zjawisko przewodnictwa objętościowego i jakie ma konsekwencje dla interpretacji topografii EEG?
  \item Jaka struktura mózgowa jest głównym generatorem dominującego rytmu alfa przy zamkniętych oczach?
  \item Podaj przykład nieprawidłowości EEG charakterystycznej dla schizofrenii oraz dla otępienia.
\end{enumerate}

\section{Literatura}
\begin{enumerate}[nosep]
  \item Szumska, I. (2019). \emph{Neurofizjologia i neuropsychologia -- skrypt dla studentów AEH}. Wydawnictwo AEH.
  \item Kalat, J. W. (2019). \emph{Biological Psychology}. Cengage Learning.
  \item Biasiucci, A., Franceschiello, B., \& Murray, M. M. (2019). Electroencephalography. \emph{Current Biology, 29}(3), R80--R85.
\end{enumerate}

\end{document}
